\section{Wind Farm Type IV: IEA37 Case Study~1}
\label{sec:iea37}
\label{sec:farm-IV}
\subsection{Experimental Setting}
The 16- and 36-turbine scenarios of IEA37 Case Study~1~\cite{Baker2019,IEA37repo} test optimizers rather than represent a site:
\begin{itemize}
\item circular sites of radii 1300 and 2000~m, $\ell_{\min}=2D$;
\item the official simplified Bastankhah Gaussian wake and wind rose (16 directions at one wind speed);
\item 6{,}030 and 30{,}030 evaluations, 30 seed-paired runs, random initialization.
\end{itemize}
In the case study the participants chose their own optimizers and budgets~\cite{Baker2019}.

% generated by analysis/rev4_farm_tables.py; regenerate, do not edit
\begin{table*}[!htb]
\centering
\caption{Type~IV farms (IEA37 Case Study~1, 16 and 36 turbines): AEP, feasibility, rank and seed-paired tests of PSO-VNS against the other eight methods.}
\label{tab:farm-IV}
\scriptsize\setlength{\tabcolsep}{1.6pt}
\begin{tabular}{lcccccccccc}
\toprule
& \multicolumn{5}{c}{6{,}030 evaluations} & \multicolumn{5}{c}{30{,}030 evaluations} \\
\cmidrule(lr){2-6}\cmidrule(lr){7-11}
Method & Feas. & Mean\,$\pm$\,SD & Best & Rank & $p_{\rm Holm}$ & Feas. & Mean\,$\pm$\,SD & Best & Rank & $p_{\rm Holm}$ \\
\midrule
\multicolumn{11}{l}{\emph{16 turbines ($r=1300$~m)}; published: example 366{,}942, best strict 418{,}924, best projected 421{,}451} \\
PSO-VNS & 30/30 & \textbf{395{,}032\,$\pm$\,4{,}587} & 402{,}782 & \textbf{1} & -- & 30/30 & \textbf{402{,}583\,$\pm$\,4{,}678} & \textbf{413{,}332} & \textbf{1} & -- \\
PSO & 30/30 & 390{,}540\,$\pm$\,7{,}192 & 403{,}753 & 3 & $6.0\times10^{-6}$$^{+}$ & 30/30 & 401{,}347\,$\pm$\,4{,}764 & 412{,}702 & 3 & 0.006$^{+}$ \\
SSA-VNS & 30/30 & 388{,}563\,$\pm$\,8{,}564 & 403{,}251 & 4 & $7.6\times10^{-4}$$^{+}$ & 30/30 & 396{,}921\,$\pm$\,6{,}267 & 406{,}925 & 4 & $4.4\times10^{-4}$$^{+}$ \\
SSA & 30/30 & 378{,}073\,$\pm$\,9{,}076 & 395{,}150 & 7 & $5.6\times10^{-8}$$^{+}$ & 30/30 & 392{,}876\,$\pm$\,4{,}394 & 401{,}459 & 7 & $3.7\times10^{-7}$$^{+}$ \\
LX-SSA & 30/30 & 373{,}977\,$\pm$\,7{,}867 & 393{,}285 & 8 & $1.5\times10^{-8}$$^{+}$ & 30/30 & 390{,}267\,$\pm$\,5{,}052 & 398{,}770 & 8 & $1.5\times10^{-8}$$^{+}$ \\
DE & 27/30 & 338{,}611\,$\pm$\,8{,}837 & 357{,}391 & 9 & $1.5\times10^{-8}$$^{+}$ & 30/30 & 359{,}952\,$\pm$\,12{,}352 & 397{,}403 & 9 & $2.6\times10^{-8}$$^{+}$ \\
GA & 30/30 & 391{,}038\,$\pm$\,5{,}425 & 399{,}691 & 2 & 0.005$^{+}$ & 30/30 & 401{,}766\,$\pm$\,3{,}288 & 409{,}194 & 2 & 0.262 \\
VNS & 30/30 & 387{,}637\,$\pm$\,5{,}591 & 400{,}169 & 5 & $4.9\times10^{-5}$$^{+}$ & 30/30 & 396{,}192\,$\pm$\,4{,}928 & 404{,}076 & 5 & $8.0\times10^{-5}$$^{+}$ \\
MS-SLSQP & 18/30 & 382{,}605\,$\pm$\,20{,}279 & \textbf{405{,}434} & 6 & $1.0\times10^{-4}$$^{+}$ & 29/30 & 394{,}567\,$\pm$\,15{,}316 & 410{,}350 & 6 & 0.010$^{+}$ \\
\midrule
\multicolumn{11}{l}{\emph{36 turbines ($r=2000$~m)}; published: example 737{,}883, best strict 863{,}676, best projected 882{,}383} \\
PSO-VNS & 30/30 & 736{,}832\,$\pm$\,19{,}268 & 787{,}709 & 3 & -- & 30/30 & 807{,}446\,$\pm$\,11{,}697 & 827{,}385 & 2 & -- \\
PSO & 27/30 & 729{,}189\,$\pm$\,21{,}150 & 787{,}989 & 4 & $6.8\times10^{-4}$$^{+}$ & 30/30 & 802{,}252\,$\pm$\,13{,}260 & 830{,}256 & 3 & $4.9\times10^{-5}$$^{+}$ \\
SSA-VNS & 30/30 & 728{,}870\,$\pm$\,22{,}859 & 766{,}803 & 5 & 0.248 & 30/30 & 791{,}404\,$\pm$\,14{,}356 & 819{,}931 & 5 & $4.2\times10^{-4}$$^{+}$ \\
SSA & 30/30 & 717{,}258\,$\pm$\,19{,}300 & 750{,}707 & 6 & $7.7\times10^{-4}$$^{+}$ & 30/30 & 771{,}298\,$\pm$\,14{,}756 & 802{,}538 & 6 & $2.6\times10^{-8}$$^{+}$ \\
LX-SSA & 30/30 & 709{,}661\,$\pm$\,19{,}845 & 746{,}301 & 7 & $4.0\times10^{-4}$$^{+}$ & 30/30 & 752{,}030\,$\pm$\,18{,}838 & 789{,}311 & 7 & $2.6\times10^{-8}$$^{+}$ \\
DE & 0/30 & -- & -- & 9 & $1.5\times10^{-8}$$^{+}$ & 0/30 & -- & -- & 9 & $1.5\times10^{-8}$$^{+}$ \\
GA & 30/30 & \textbf{774{,}972\,$\pm$\,8{,}083} & \textbf{795{,}530} & \textbf{1} & $3.9\times10^{-8}$$^{-}$ & 30/30 & \textbf{823{,}061\,$\pm$\,6{,}994} & \textbf{836{,}526} & \textbf{1} & $1.9\times10^{-6}$$^{-}$ \\
VNS & 30/30 & 741{,}566\,$\pm$\,24{,}100 & 785{,}791 & 2 & 0.404 & 30/30 & 801{,}797\,$\pm$\,11{,}876 & 824{,}230 & 4 & 0.119 \\
MS-SLSQP & 4/30 & 733{,}342\,$\pm$\,13{,}798 & 748{,}571 & 8 & $1.1\times10^{-7}$$^{+}$ & 14/30 & 798{,}913\,$\pm$\,41{,}511 & 833{,}365 & 8 & $2.4\times10^{-4}$$^{+}$ \\
\bottomrule
\end{tabular}
\par\vspace{2pt}\parbox{\textwidth}{\scriptsize Random initialization, 30 seed-paired runs (seeds 1--30) per method and budget; stored feasibility labels. Feas.: feasible runs; Mean\,$\pm$\,SD: mean AEP (MWh, rounded to integers) and its standard deviation over the feasible runs (mean AEP losses relative to the wake-free AEP: rev4\_farm\_tables.json); Best: best feasible run; --: no feasible run. Rank: feasibility-aware rank of the mean AEP within the methods of the pool (ranking rule of Section~\ref{sec:setup-stats}: a method with fewer than 15 feasible runs ranks below the qualified methods, by its number of feasible runs). Bold: rank-1 method (mean, rank), highest best run. $p_{\rm Holm}$: two-sided seed-paired Wilcoxon signed-rank test of PSO-VNS against the method (feasibility first: every infeasible run below every feasible run, infeasible runs ordered by minimum spacing), Holm-adjusted over the comparisons of PSO-VNS at that budget (8 comparisons); $^{+}$/$^{-}$: PSO-VNS significantly better/worse ($\alpha=0.05$). Official IEA37 Gaussian-wake model and wind rose (16 directions, one wind speed), circular boundary, minimum spacing $2D$. Published layouts~\cite{Baker2019,IEA37repo}: example layout, best submitted layout with every turbine within 1~mm of the boundary (strict) and after radial projection onto it. Run-level Friedman tests over the seeds: 16 turbines, 6{,}030: $\chi^2_F=126.8$ (8 d.f.), $p=1.3\times10^{-23}$; 16 turbines, 30{,}030: $\chi^2_F=132.5$ (8 d.f.), $p=8.4\times10^{-25}$; 36 turbines, 6{,}030: $\chi^2_F=170.0$ (8 d.f.), $p=1.3\times10^{-32}$; 36 turbines, 30{,}030: $\chi^2_F=158.7$ (8 d.f.), $p=3.0\times10^{-30}$.}
\end{table*}

\subsection{Numerical Results}
With 16 turbines PSO-VNS ranks first at both budgets (Table~\ref{tab:farm-IV}). With 36 turbines GA ranks first at both budgets, and PSO-VNS ranks third at 6,030 and second at 30,030. With forward differences, every MS-SLSQP gradient costs $2N+1$ evaluations, so 6,030 evaluations allow at most 182 gradients with 16 turbines and 82 with 36, and MS-SLSQP is feasible in only 4 and 14 of 30 runs with 36 turbines. None of the methods reaches the best feasible published layouts (Table~\ref{tab:farm-IV}; all gaps: Table~\ref{S-tab:F-iea-gap}); as the participants' budgets are not reported comparably, the size of the gap is indicative only.
% CHECK-FINAL [C40]: best of all methods below the best feasible published layout (strict) in both scenarios

\subsection{Statistical Analysis}
With 16 turbines PSO-VNS is better than all eight other methods at 6,030 evaluations and than all but GA at 30,030 ($p_{\rm Holm}=0.262$). With 36 turbines GA is better than PSO-VNS ($p_{\rm Holm}=3.9\times10^{-8}$ and $1.9\times10^{-6}$); PSO-VNS does not differ significantly from VNS at either budget.

On this analytic model SLSQP can also use the exact gradient of the IEA37 AEP. With it, as multistart MS-SLSQP and after the PSO phase of PSO-VNS (PSO-SLSQP), an exact-gradient variant ranks first in both scenarios and at both budgets and beats every method without exact gradients in all 30 seed pairs (Section~\ref{S-sec:S-xp-grad}, Table~\ref{S-tab:S-sa-iea-pool}), a result that rests on the feasibility labels reported by these runs, which the rounded stored coordinates cannot decide (Section~\ref{sec:archive-reliability}).

\subsection{Exact Gradients on IEA37}
\label{sec:iea-grad}
With forward differences, every MS-SLSQP gradient costs $2N$ evaluations besides the objective value, $2N+1$ in all, so 6,030 evaluations allow at most 182 gradients with 16 turbines and 82 with 36; exact gradients remove most of this cost, which grows with $N$ (Section~\ref{sec:related-wflop}). On the IEA37 farms (Type~IV; Section~\ref{sec:iea37}), whose model has an analytic AEP gradient, we therefore also ran SLSQP with exact objective and constraint derivatives, as multistart (exact-gradient MS-SLSQP) and after the PSO phase of PSO-VNS (PSO-SLSQP), each gradient charged $c_g=3$ evaluations, its measured cost (30 seeds, random starts; Section~\ref{S-sec:S-xp-grad}, Tables~\ref{S-tab:S-xp-grad} and~\ref{S-tab:S-xp-grad-calls}).

In a joint pool of 13 methods (the nine methods, LX-SSA-VNS, RS-VNS and the two exact-gradient variants; Table~\ref{S-tab:S-sa-iea-pool}), an exact-gradient method ranks first in every scenario and budget and is better than each of the 11 other methods in all 30 seed pairs ($p_{\rm Holm}=2.2\times10^{-8}$); the two exact-gradient variants differ significantly only at 30,030 evaluations. With 36 turbines, their mean AEP at 6,030 evaluations (832{,}893.3 and 832{,}404.0~MWh) exceeds that of every other method at 30,030 (highest: GA, 823{,}061.5~MWh), and the best PSO-SLSQP layouts at 30,030 evaluations lie 0.04\% (16 turbines) and 1.84\% (36 turbines) below the best feasible published layouts (strict convention). Three qualifications apply:
\begin{enumerate}
\item The comparison tests complete implementations of the second phase, not the isolated effect of the AEP derivative: the variants also differentiate the constraints (MS-SLSQP of the main comparison uses finite differences), PSO-SLSQP also restarts from the incumbent and personal bests, and the objective is only piecewise smooth (the derivative treats the upstream-wake indicator as locally constant).
\item The budget is equal in charged objective-equivalent units, not in time: SLSQP subproblems, constraint evaluations and their Jacobians are not charged, and per run the exact-gradient methods take 2.3--5.6 times as long as PSO-VNS (medians; Table~\ref{S-tab:S-sa-time-sites}).
\item The results rest on the feasibility labels reported by the runs. The constraints are active at the optimum, so the rounded stored coordinates cannot decide the strict labels: full-precision reruns resolved 20 of the 240 labels (all confirmed) and left 220 undecidable (Section~\ref{sec:archive-reliability}). For these 220 runs the spacing constraint is satisfied by the minimum spacing stored at full precision, the rounding bounds any boundary violation by 1.34~mm ($\le1.0\times10^{-6}$ of the radius), and no label is decided infeasible at any tolerance (Table~\ref{S-tab:S-sa-feasbounds}). If every open label were infeasible, the exact-gradient methods would keep at most 7 feasible runs of 30 per setting and PSO-VNS (16 turbines) or GA (36 turbines) would lead; with at least 15 of their 30 runs feasible they keep ranks 1--2 in every setting, although no longer significantly ahead of most other methods (Table~\ref{S-tab:S-sa-feasworst}).
\end{enumerate}
